\chapter{حاسوب من خلايا حية}
% full version: chapters/21_living.tex

\pencilfig{21}{\pencilart{21}}{عصبونات حية على رقاقة: تحتها آلاف الأقطاب، وكل منها يستطيع أن يسمع وأن يتكلم.}

هل يمكن بناء آلة لا من السيليكون، بل من عصبونات حية؟ يمكن، وهي تُبنى فعلًا. تُستنبت العصبونات مباشرة على رقاقة فيها آلاف الأقطاب، طبقةً مسطحة أو كتلًا صغيرة تشبه دماغًا ضئيلًا. كل قطب يستطيع أن يسمع وأن يرسل إشارة. إنها منصة اختبار بدائرة مكشوفة: كل شيء يمكن قياسه، وفي كل شيء يمكن التدخل. لكن لا محطات إذاعية فيها: إنها ناقل بيانات بلا سجلات تحكم.

\section*{ما تفعله الشبكة وحدها}

إذا تُركت المزرعة لنفسها، اشتعلت كلها من حين لآخر: تنقر العصبونات كلها تقريبًا معًا، ثم تهدأ، بفواصل من ثانية إلى دقائق. وهذا مذبذب مألوف: الشبكة تثير نفسها وتتعب، إلى أن تستعيد الوصلات مخزونها من المادة.

\pencilfig{129}{%
% spike raster of 8 neurons in a culture left to itself: sparse single clicks and shared bursts
\foreach \b in {0.9,3.1,4.0,6.6}{\fill[pcAmber, opacity=0.18] (\b-0.1,-0.15) rectangle (\b+0.32,2.95);}
\foreach \y/\ss in {0/{0.96,1.0,1.13,2.43,3.0,3.12,3.24,4.02,4.09,4.15,6.66,6.69,6.81},
  1/{0.46,0.88,1.0,1.05,3.09,3.2,3.94,4.04,4.37,6.65,6.71},
  2/{0.73,0.84,0.94,1.41,3.08,3.12,3.21,4.05,4.06,4.17,6.59,6.63},
  3/{0.89,0.93,1.06,3.1,3.19,3.28,3.89,3.94,4.13,4.2,4.31,6.63,6.65,6.78},
  4/{0.87,0.96,3.09,3.15,3.23,3.73,3.96,4.06,4.19,5.98,6.66,6.68,6.75},
  5/{0.44,0.92,0.97,3.05,3.22,3.27,4.01,4.07,4.17,5.76,6.57,6.73,6.75},
  6/{0.92,1.0,1.73,2.85,3.18,3.19,4.0,4.07,6.53,6.66,6.73},
  7/{0.87,0.92,3.13,3.13,3.27,3.99,4.06,4.17,6.44,6.61,6.72,7.13}}{
  \foreach \s in \ss {\draw[pcBlue, line width=0.7pt] (\s,0.35*\y) -- (\s,0.35*\y+0.25);}}
\draw[arr] (0,-0.3) -- (7.5,-0.3); \node[lbl, anchor=north] at (3.75,-0.32) {الزمن};
\node[lbl, anchor=east, align=right] at (-0.05,1.35) {ثمانية\\عصبونات};
\node[lbl, text=pcAmber!70!black, anchor=south] at (3.55,2.95) {رشقات شبكية};
}{مزرعة على رقاقة متروكة لنفسها. النقرات المنفردة نادرة، ومن حين لآخر تشتعل الشبكة كلها تقريبًا دفعة واحدة، ثم تهدأ إلى أن تستعيد الوصلات مخزونها.}

\section*{معلم بلا دوبامين}

كيف نعلّم شبكة كهذه إذا لم يكن فيها دوبامين؟ كانت الطرق الأولى كهربائية. مثلًا: تحفيز الشبكة إلى أن تعطي الاستجابة المطلوبة، ثم التوقف فورًا. فتترسخ الاستجابة. وفي تجربة أخرى كانت المزرعة تتحكم في ”جسم“ افتراضي بسيط وتتعلم قيادته إلى الهدف. وفي كل هذه التجارب المعلم نمط من التيار يختاره المهندس.

\section*{دوبامين بومضة ضوء}

ذات مرة حاولوا جعل المعلم كيميائيًا. على إحدى المنصات ذات الكتل العصبية الحية أضيف إلى السائل المغذي دوبامين داخل ”قفص“ حساس للضوء يمنعه من العمل. وحين كانت الشبكة تستجيب للمنبه أقوى من قبل، كانت ”تُكافأ“: ومضة من الأشعة فوق البنفسجية تحطم القفص، فيخرج الدوبامين. وخلال مئتين وأربعين تكرارًا نمت الاستجابات إجمالًا. لكن المؤلفين أنفسهم يكتبون بصدق أنه لا يمكن بتجربة واحدة كهذه الجزم بأن الدوبامين هو السبب.

\section*{عمود على عربة}

في تجربة حديثة عُلّمت كتل عصبية موازنة عمود على عربة، وهي مسألة كلاسيكية لبرامج التعلم. وكان برنامج يختار نبضات التعليم التي تُعطى. عملت الإشارات التي اختارها البرنامج أفضل من العشوائية، لكن ما تُعلّم كان يختفي بعد ثلاثة أرباع ساعة من الراحة. وإذا حُجبت وصلات \nt{GLUT}، لم يحدث تعلم أصلًا: ما تُعلّم يسكن هناك بالذات. أما المعلم فلا يزال في الخارج، لكنه الآن برنامج لا مهندس.

\section*{الخزان الحي}

وثمة نهج آخر: لا تُعلَّم الشبكة الحية أصلًا. تُستخدم ”خزانًا“: نظامًا غنيًا متشابكًا، يتعلم عدّاد إلكتروني بسيط تفكيك استجاباته. هكذا ساعدت كتلة حية على تمييز الأصوات. لكن المؤلفين يقولون إن الوصلات في الكتلة نفسها كانت تتغير هي الأخرى في أثناء العمل.

\section*{الثمن والسؤال}

مليون عصبون على رقاقة يستهلك نحو جزء من عشرة آلاف من الواط، أقل بكثير من الإلكترونيات المحيطة بها. والمنصة لا تستطيع وحدها أن تقرر ما يُعدّ نجاحًا: هذا الدور يؤديه دائمًا أحد من الخارج. ومع نمو هذه الأنظمة يُطرح سؤال أخلاقي لم يحسمه أحد حتى الآن.

\fullversion{21}
